% =====================================================================
\chapter{Probabilidad, condicionamiento y teorema de Bayes}
\label{tema:1}
% =====================================================================

Este tema fija el lenguaje del módulo: espacios de probabilidad, probabilidad condicionada, independencia y teorema de Bayes. La presentación sigue el capítulo~1 de~\cite{bt} y las lecciones~1 a~3 de~\cite{mit6041}; la terminología bayesiana (a priori, verosimilitud, a posteriori) y la forma de razón de probabilidades siguen las clases~3 y~12 de~\cite{mit1805} y el capítulo~2 de~\cite{bh}. Deisenroth \emph{et al.}~\cite[§6.1 y §6.3]{mml} ofrecen el mismo material con orientación al aprendizaje automático.

\section{Modelos probabilísticos}

\begin{definicion}[Espacio de probabilidad]\label{def:espacio}
Un \emph{espacio de probabilidad} es una terna $(\Omega,\mathcal{F},\PP)$ formada por un \emph{espacio muestral} $\Omega$ (el conjunto de resultados posibles de un experimento), una familia $\mathcal{F}$ de subconjuntos de $\Omega$, llamados \emph{sucesos}, cerrada por complementos y uniones numerables y que contiene a $\Omega$, y una función $\PP\colon\mathcal{F}\to\R$ que cumple los \emph{axiomas de Kolmogorov}:
\begin{enumerate}
  \item \emph{no negatividad}: $\PP(A)\ge 0$ para todo $A\in\mathcal{F}$;
  \item \emph{normalización}: $\PP(\Omega)=1$;
  \item \emph{aditividad numerable}: si $A_1,A_2,\dots$ son disjuntos dos a dos, $\PP\bigl(\bigcup_{i}A_i\bigr)=\sum_{i}\PP(A_i)$.
\end{enumerate}
\end{definicion}

Cuando $\Omega$ es finito o numerable se toma $\mathcal{F}$ igual a todos los subconjuntos de $\Omega$, y basta dar la probabilidad de cada resultado: $\PP(A)=\sum_{\omega\in A}\PP(\{\omega\})$. Si además $\Omega$ es finito y todos los resultados son equiprobables, $\PP(A)=\abs{A}/\abs{\Omega}$ (ley de Laplace)~\cite[§1.2]{bt}.

\begin{proposicion}[Consecuencias de los axiomas]\label{prop:axiomas}
Para sucesos $A,B\in\mathcal{F}$:
\begin{enumerate}
  \item $\PP(\varnothing)=0$ y $\PP(A^{c})=1-\PP(A)$;
  \item si $A\subseteq B$, entonces $\PP(A)\le\PP(B)$;
  \item $\PP(A\cup B)=\PP(A)+\PP(B)-\PP(A\cap B)$;
  \item $\PP(A\cup B)\le\PP(A)+\PP(B)$ (cota de la unión).
\end{enumerate}
\end{proposicion}

\begin{proof}
(i) $\Omega$ y $\varnothing$ son disjuntos y $\Omega\cup\varnothing=\Omega$; por aditividad, $\PP(\Omega)=\PP(\Omega)+\PP(\varnothing)$, luego $\PP(\varnothing)=0$. Como $A$ y $A^{c}$ son disjuntos y su unión es $\Omega$, $\PP(A)+\PP(A^{c})=1$. (ii) $B=A\cup(B\cap A^{c})$ con unión disjunta, así que $\PP(B)=\PP(A)+\PP(B\cap A^{c})\ge\PP(A)$ por no negatividad. (iii) $A\cup B=A\cup(B\cap A^{c})$ y $B=(A\cap B)\cup(B\cap A^{c})$, ambas uniones disjuntas; restando las dos igualdades de probabilidades se obtiene el resultado. (iv) Se sigue de (iii), porque $\PP(A\cap B)\ge 0$.
\end{proof}

\section{Probabilidad condicionada}

\begin{definicion}[Probabilidad condicionada]
Sean $A,B\in\mathcal{F}$ con $\PP(A)>0$. La \emph{probabilidad de $B$ condicionada a $A$} es
\[
\PP(B\mid A)=\frac{\PP(A\cap B)}{\PP(A)}.
\]
\end{definicion}

Condicionar a $A$ equivale a tomar $A$ como nuevo espacio muestral y renormalizar. Para $A$ fijo, $B\mapsto\PP(B\mid A)$ cumple los tres axiomas de Kolmogorov~\cite[§1.3]{bt}; en consecuencia, todas las propiedades de la \cref{prop:axiomas} valen también para probabilidades condicionadas. Despejando se obtiene la \emph{regla del producto}, $\PP(A\cap B)=\PP(A)\,\PP(B\mid A)$, cuya iteración es la regla de la cadena del módulo~1.

\begin{observacion}
Un modelo de lenguaje es una familia de probabilidades condicionadas $\PP(x_t\mid x_{1:t-1})$: el suceso condicionante es el contexto observado y la distribución condicionada es la del token siguiente. Todo lo que se demuestre en este tema para probabilidades condicionadas se aplica, por tanto, a las salidas de un LLM.
\end{observacion}

\section{Probabilidad total y regla de Bayes}

\begin{definicion}[Partición]
Una familia finita o numerable de sucesos $H_1,H_2,\dots$ es una \emph{partición} de $\Omega$ si son disjuntos dos a dos y su unión es $\Omega$.
\end{definicion}

\begin{teorema}[Probabilidad total]\label{teo:total}
Sea $H_1,H_2,\dots$ una partición de $\Omega$ con $\PP(H_i)>0$ para todo $i$. Para todo suceso $D$,
\[
\PP(D)=\sum_{i}\PP(D\mid H_i)\,\PP(H_i).
\]
\end{teorema}

\begin{proof}
Los sucesos $D\cap H_i$ son disjuntos dos a dos y su unión es $D\cap\Omega=D$. Por aditividad, $\PP(D)=\sum_i\PP(D\cap H_i)$, y por la regla del producto, $\PP(D\cap H_i)=\PP(H_i)\,\PP(D\mid H_i)$.
\end{proof}

\begin{teorema}[Regla de Bayes]\label{teo:bayes}
Con las hipótesis del \cref{teo:total} y $\PP(D)>0$, para cada $j$,
\[
\PP(H_j\mid D)=\frac{\PP(D\mid H_j)\,\PP(H_j)}{\sum_{i}\PP(D\mid H_i)\,\PP(H_i)}.
\]
\end{teorema}

\begin{proof}
Por la definición y la regla del producto aplicada de las dos formas posibles, $\PP(H_j\mid D)=\PP(H_j\cap D)/\PP(D)=\PP(D\mid H_j)\PP(H_j)/\PP(D)$. Basta sustituir $\PP(D)$ por la expresión del \cref{teo:total}.
\end{proof}

En el lenguaje de la inferencia, los $H_i$ son \emph{hipótesis} y $D$ son los \emph{datos}. $\PP(H_j)$ es la probabilidad \emph{a priori}, $\PP(D\mid H_j)$ la \emph{verosimilitud}, $\PP(H_j\mid D)$ la probabilidad \emph{a posteriori} y el denominador, $\PP(D)$, la \emph{evidencia}, que solo normaliza. Abreviadamente, la probabilidad a posteriori es proporcional al producto de verosimilitud y probabilidad a priori~\cite[clase~11]{mit1805}.

\begin{ejemplo}[Filtro de correo no deseado]\label{ej:spam}
El $20\,\%$ del correo es no deseado ($S$). La palabra «gratis» ($G$) aparece en el $50\,\%$ de esos mensajes y en el $5\,\%$ de los demás. Por el \cref{teo:bayes},
\[
\PP(S\mid G)=\frac{0{,}5\cdot 0{,}2}{0{,}5\cdot 0{,}2+0{,}05\cdot 0{,}8}=\frac{0{,}10}{0{,}14}=\frac57\approx 0{,}714.
\]
La \cref{fig:arbol} representa el cálculo como un árbol: cada camino lleva el producto de sus probabilidades, y el resultado es la fracción de la masa de los caminos que terminan en $G$ que corresponde al camino que pasa por $S$.
\end{ejemplo}

\begin{figure}[h]
\centering
\begin{tikzpicture}[font=\small,level distance=3.2cm,
  level 1/.style={sibling distance=2.6cm},level 2/.style={sibling distance=1.25cm},
  grow=right,edge from parent/.style={draw,-{Stealth[length=2mm]}}]
\node[draw,circle,inner sep=1.5pt] {}
  child {node[right] {$S^{c}$}
    child {node[right] {$G^{c}$: \ $0{,}8\cdot 0{,}95=0{,}76$} edge from parent node[below,pos=0.6] {$0{,}95$}}
    child {node[right] {$G$: \ $0{,}8\cdot 0{,}05=0{,}04$} edge from parent node[above,pos=0.6] {$0{,}05$}}
    edge from parent node[below left] {$0{,}8$}}
  child {node[right] {$S$}
    child {node[right] {$G^{c}$: \ $0{,}2\cdot 0{,}5=0{,}10$} edge from parent node[below,pos=0.6] {$0{,}5$}}
    child {node[right] {$G$: \ $0{,}2\cdot 0{,}5=0{,}10$} edge from parent node[above,pos=0.6] {$0{,}5$}}
    edge from parent node[above left] {$0{,}2$}};
\end{tikzpicture}
\caption{Árbol de probabilidades del \cref{ej:spam}. La evidencia es $\PP(G)=0{,}10+0{,}04=0{,}14$ y la probabilidad a posteriori, $0{,}10/0{,}14\approx 0{,}714$.}
\label{fig:arbol}
\end{figure}

Cuando solo hay dos hipótesis, $H$ y $H^{c}$, es cómodo trabajar con la \emph{razón de probabilidades} (en inglés, \emph{odds}) $O(H)=\PP(H)/\PP(H^{c})$.

\begin{proposicion}[Bayes en forma de razón]\label{prop:odds}
Si $\PP(D\mid H^{c})>0$, entonces
\[
O(H\mid D)=\frac{\PP(D\mid H)}{\PP(D\mid H^{c})}\cdot O(H),
\]
donde $O(H\mid D)=\PP(H\mid D)/\PP(H^{c}\mid D)$. El cociente $\PP(D\mid H)/\PP(D\mid H^{c})$ se llama \emph{factor de Bayes} o razón de verosimilitudes.
\end{proposicion}

\begin{proof}
Se aplica el \cref{teo:bayes} a $H$ y a $H^{c}$ y se dividen las dos expresiones: la evidencia $\PP(D)$ se cancela y queda $\dfrac{\PP(H\mid D)}{\PP(H^{c}\mid D)}=\dfrac{\PP(D\mid H)\,\PP(H)}{\PP(D\mid H^{c})\,\PP(H^{c})}$.
\end{proof}

La forma de razón muestra con claridad la llamada \emph{falacia de la probabilidad a priori}: un factor de Bayes grande no basta para que la hipótesis sea probable si $O(H)$ es muy pequeña (\cref{eje:1-detector,eje:1-cribado}).

\section{Independencia}

\begin{definicion}[Independencia]\label{def:indep}
Dos sucesos $A$ y $B$ son \emph{independientes} si $\PP(A\cap B)=\PP(A)\,\PP(B)$. Si $\PP(A)>0$, la condición equivale a $\PP(B\mid A)=\PP(B)$: saber que ha ocurrido $A$ no cambia la probabilidad de $B$. Los sucesos $A_1,\dots,A_n$ son \emph{independientes} si $\PP\bigl(\bigcap_{i\in S}A_i\bigr)=\prod_{i\in S}\PP(A_i)$ para todo subconjunto $S\subseteq\{1,\dots,n\}$.
\end{definicion}

\begin{definicion}[Independencia condicionada]
Sea $C$ con $\PP(C)>0$. Los sucesos $A$ y $B$ son \emph{independientes condicionalmente a $C$} si $\PP(A\cap B\mid C)=\PP(A\mid C)\,\PP(B\mid C)$.
\end{definicion}

\begin{observacion}
La independencia condicionada no implica la independencia, ni al revés. Dos lanzamientos de una moneda elegida al azar entre una equilibrada y una trucada son independientes condicionalmente a la moneda, pero no lo son incondicionalmente: una cara en el primer lanzamiento hace más probable la moneda que favorece la cara y, con ella, una cara en el segundo (\cref{eje:1-monedas}). Del mismo modo, independencia no es lo mismo que incompatibilidad: dos sucesos disjuntos con probabilidad positiva nunca son independientes (\cref{eje:1-disjuntos}).
\end{observacion}

\begin{observacion}[Conexión con los LLM]
El aprendizaje en contexto, es decir, la capacidad de un LLM de adaptar su respuesta a ejemplos incluidos en la instrucción sin modificar sus pesos (módulo~22), se ha interpretado como una inferencia bayesiana implícita: el contexto actúa como los datos $D$ y el modelo concentra su distribución sobre las continuaciones compatibles con un «concepto» latente $H$~\cite{xie}. Se trata de un modelo teórico explicativo, no de una descripción del funcionamiento interno de un modelo comercial concreto.
\end{observacion}

\begin{lecturas}
\item Bertsekas y Tsitsiklis~\cite{bt}, §1.2 «Probabilistic Models», §1.3 «Conditional Probability», §1.4 «Total Probability Theorem and Bayes' Rule» y §1.5 «Independence».
\item MIT 6.041SC~\cite{mit6041}, unidad~I, lecciones~1 «Probability Models and Axioms», 2 «Conditioning and Bayes' Rule» y 3 «Independence».
\item MIT 18.05~\cite{mit1805}, clase~3 «Conditional Probability, Independence and Bayes' Theorem» y clase~12b «Bayesian Updating: Odds».
\item Blitzstein y Hwang~\cite{bh}, cap.~2 «Conditional Probability».
\item Deisenroth, Faisal y Ong~\cite{mml}, §6.1 «Construction of a Probability Space» y §6.3 «Sum Rule, Product Rule, and Bayes' Theorem».
\end{lecturas}

\section*{Ejercicios del tema 1}
\addcontentsline{toc}{section}{Ejercicios del tema 1}

\begin{ejercicio}{1}{eje:1-inversa}
Proponga un ejemplo cotidiano en el que $\PP(A\mid B)$ y $\PP(B\mid A)$ sean muy distintas. Asigne valores plausibles a $\PP(A)$, $\PP(B)$ y $\PP(B\mid A)$, calcule $\PP(A\mid B)$ e identifique qué cociente relaciona ambas probabilidades condicionadas.
\end{ejercicio}

\begin{ejercicio}{1}{eje:1-disjuntos}
Demuestre que si $A$ y $B$ son disjuntos y $\PP(A)>0$, $\PP(B)>0$, entonces $A$ y $B$ no son independientes. Explique intuitivamente el resultado.
\end{ejercicio}

\begin{ejercicio}{2}{eje:1-detector}
Un detector de texto generado por IA tiene una sensibilidad del $99\,\%$ (marca el $99\,\%$ de los textos generados) y una tasa de falsos positivos del $5\,\%$. Si el $1\,\%$ de los textos analizados son generados, ¿qué probabilidad hay de que un texto marcado lo sea realmente?
\end{ejercicio}

\begin{ejercicio}{2}{eje:1-cribado}
Una enfermedad afecta al $0{,}5\,\%$ de la población. Una prueba diagnóstica da un $5\,\%$ de falsos positivos y un $2\,\%$ de falsos negativos. Para una persona con resultado positivo, calcule la razón de probabilidades a priori, el factor de Bayes, la razón a posteriori y la probabilidad a posteriori de estar enferma.
\fuente{adaptado de MIT 18.05~\cite{mit1805}, clase~12, problema~2.}
\end{ejercicio}

\begin{ejercicio}{2}{eje:1-monedas}
En una bolsa hay dos monedas de aspecto idéntico: una equilibrada y otra con probabilidad de cara $1/4$. Se extrae una al azar y se lanza dos veces. (a) Si salen dos caras, ¿qué probabilidad hay de que sea la moneda equilibrada? (b) ¿Son independientes los sucesos «cara en el primer lanzamiento» y «cara en el segundo»? Justifíquelo con un cálculo.
\fuente{adaptado de Harvard Stat 110~\cite{stat110}, \emph{Homework}~2, problema~4.}
\end{ejercicio}

\begin{ejercicio}{3}{eje:1-complementos}
Demuestre que si $A$ y $B$ son independientes, también lo son $A$ y $B^{c}$, y $A^{c}$ y $B^{c}$.
\fuente{adaptado de Harvard Stat 110~\cite{stat110}, \emph{Strategic Practice}~2, sección «Independence», problema~2.}
\end{ejercicio}
